\documentclass[10pt,a4]{examenEPN0}
\usepackage{comandosEPN}

\usepackage{empheq}

\newcommand{\itemcp}{\stepcounter{enumi}\item[ \textcolor{red}{\theenumi.$^*$}] }

\newcommand{\iitemcp}{\stepcounter{enumii}\item[ \textcolor{red}{\theenumii)$^*$}] }


\materia{Álgebra Lineal}
%\autor{Departamento de Formación Básica}
\autor{GR \hspace{1cm}.}
%\fecha{Semestre 2026-A}
\fecha{Nombre: \hspace{9cm} Cédula:}
\logouno[1.6cm]{Logo_EPN.eps}
\logodos[2.6cm]{Logo_DFB.eps}

\newcommand{\la}{\lambda}

%Datos
\newcommand{\datosAlgebra}[2]{
\EnExamen
{
\vspace{5mm}
    \noindent
    Nombres: \hrulefill\
    Carrera: \rule{4.5cm}{0.4pt}\
    \ifthenelse{\equal{#1}{}}%
    {Paralelo: \underline{\parbox{1.8cm}{\parbox{1cm}{}}}\\[3mm]}
    {Paralelo: \underline{\parbox{1.8cm}{\centering #1}}\\[3mm]}
    C\'edula: \rule{5cm}{0.4pt}\
    Firma: \hrulefill\
    Fecha: \underline{\parbox{2.1cm}{\centering #2}}
}}


% \newcommand{\indicacionesPrueba}{%
% \EnExamen{
% \begin{indicaciones}
%     \item La prueba tiene una duraci\'on  de \textbf{120 minutos.}
%     \item No se permite la utilizaci\'on de dispositivos electr\'onicos (celular, tablets, calculadora, etc.)
%     \item Leer detenidamente los enunciados, \textbf{justificar la respuesta para acceder al puntaje total}.
%     \item La prueba es completamente individual. Los estudiantes deberán dejar sus teléfonos celulares guardados y no podrán tenerlos consigo (por ejemplo, en los bolsillos de su ropa) durante toda la evaluación. En caso de detectarse algún acto de deshonestidad académica o el incumplimiento de esta norma, no podrán seguir rindiendo la prueba y su calificación será \textbf{cero}.
% \end{indicaciones}
% \pagestyle{empty}
% }}

\newcommand{\indicacionesPrueba}{%
\EnExamen{
\begin{indicaciones}
    \item La prueba tiene una duraci\'on  de \textbf{60 minutos.}
    \item Leer detenidamente los enunciados, \textbf{justificar la respuesta} para acceder al puntaje total.
    \item No se permite la utilizaci\'on de dispositivos electr\'onicos (celular, tablets, calculadora, etc.)
    \item La prueba es completamente individual
    \item No se permite portar teléfonos celulares, relojes inteligentes ni otros dispositivos electrónicos durante el examen, aun cuando permanezcan en los bolsillos. Estos dispositivos deberán guardarse en la mochila antes del inicio de la evaluación y esta deberá mantenerse cerrada durante todo el examen.
    \item He leído las indicaciones y acepto que el incumplimiento de estas se penaliza con \textbf{cero}. Firma:...................................

    
\end{indicaciones}
\pagestyle{empty}
}}

\newcommand{\indicacionesExamen}{
\EnExamen{\begin{indicaciones}
    \item La evaluación tiene una duraci\'on  de \textbf{120 minutos.}
    \item No se permite la utilizaci\'on de dispositivos electr\'onicos (celular, tablets, calculadora, etc.)
    %\item Leer detenidamente los enunciados, justificar la respuesta para acceder al puntaje total.
    \item El examen es completamente individual. En caso de detectarse alg\'un acto de deshonestidad acad\'emica no podr\'a seguir rindiendo la prueba y su calificaci\'on ser\'a \textbf{cero}.
\end{indicaciones}
}\pagestyle{empty}
}

%% Imprimir las una hoja A4, dos veces en A5
\newcommand{\imprimirDoble}
{%
\usepackage{pgfmorepages}
\pgfmorepagesloadextralayouts
\pgfpagesuselayout{repeated 2-up}[a4paper,landscape]
\thispagestyle{empty}
}


%Este sí se puede comentar
\newenvironment{solucion}{\color{blue}\proof[Solución]}{\endproof\color{black}}




% \renewcommand{\@autor}{
% ---\\
% \phantom{.} 
% ...}
% \makeatother

\examen{Examen 01: solucionario y rúbrica}
\makeatletter
\renewcommand{\@autor}{Departamento de Formación Básica}
\renewcommand{\@fecha}{GR5E -- Jueves de 7 a 9 -- 19/11/2026}
\makeatother
\begin{document}
\encabezado
\noindent\textbf{Criterio general (20 puntos).} Se asigna crédito parcial por
los pasos correctos. Un error aritmético se penaliza donde se origina;
los pasos posteriores se valoran por su coherencia, sin repetir la penalización.
Se aceptan procedimientos equivalentes debidamente justificados.

\section*{1. Alquiler de equipos (4 puntos)}
\noindent\textbf{Rúbrica.}
\begin{itemize}[itemsep=2pt]
    \item Definir las variables y unidades: 0,25 puntos; plantear las ecuaciones: 0,75 puntos (0,25 por ecuación); identificar $A$, $\mathbf{x}$ y $b$: 1 punto.
    \item Resolver con procedimiento: 1,5 puntos; interpretar y comprobar las tarifas: 0,5 puntos.
\end{itemize}
Sean $p,t,s$ las tarifas diarias en dólares de proyector, pantalla y sonido.
\[
    p+t+s=65,\qquad 2p+t=80,\qquad t-s=5.
\]
\[
    A=\begin{pmatrix}1&1&1\\2&1&0\\0&1&-1\end{pmatrix},\quad
    \mathbf{x}=\begin{pmatrix}p\\t\\s\end{pmatrix},\quad
    b=\begin{pmatrix}65\\80\\5\end{pmatrix}.
\]
De $t=s+5$ resulta $p+2s=60$ y $2p+s=75$; por tanto,
$s=15$, $p=30$ y $t=20$.
Las tarifas son \textbf{\$30, \$20 y \$15 diarios}, respectivamente.
Se verifica $30+20+15=65$, $60+20=80$ y $20-15=5$.

\section*{2. Sistema con parámetro (4 puntos)}
\noindent\textbf{Rúbrica.}
\begin{itemize}[itemsep=2pt]
    \item Calcular el determinante con procedimiento: 1,5 puntos; establecer cuándo hay solución única: 0,5 puntos.
    \item Sustituir $\alpha=2$ y reducir: 1 punto; exhibir una contradicción: 0,5 puntos; concluir que el conjunto solución es vacío: 0,5 puntos.
\end{itemize}
Sumando la primera fila a la segunda en la matriz de coeficientes,
\[
    \det(A)=\det\begin{pmatrix}1&\alpha&2\\0&\alpha-2&\alpha\\0&\alpha-2&2\alpha+2\end{pmatrix}
    =(\alpha-2)(\alpha+2).
\]
Hay solución única si y solo si $\alpha\notin\{-2,2\}$.
Para $\alpha=2$, sumar las dos primeras ecuaciones da $2x_3=0$,
mientras que la tercera exige $6x_3=1$. Son incompatibles:
\textbf{no hay solución}, y el conjunto solución es $\varnothing$.

\section*{3. Intersección de subespacios de matrices (4 puntos)}
\noindent\textbf{Rúbrica.}
\begin{itemize}[itemsep=2pt]
    \item Imponer todas las condiciones de $S$ y $T$: 1 punto; parametrizar la intersección: 1 punto.
    \item Plantear las condiciones sobre $k$: 1 punto; hallar el valor y verificar todas las condiciones: 1 punto.
\end{itemize}
Si $a_{11}=t$, las restricciones implican $a_{12}=a_{21}=t$ y
$a_{22}=-t$. Por tanto,
\[
    \boxed{S\cap T=\left\{t\begin{pmatrix}1&1\\1&-1\end{pmatrix}:t\in\R\right\}.}
\]
Para la matriz dada se exige $3=2k-1$, $2k-4=0$ y $k+1-3=0$.
Todas se cumplen exactamente cuando $k=2$, que produce
$\begin{pmatrix}3&3\\3&-3\end{pmatrix}$.

\section*{4. Espacio generado (4 puntos)}
\noindent\textbf{Rúbrica.}
\begin{itemize}[itemsep=2pt]
    \item Describir el espacio generado en forma paramétrica o cartesiana: 1 punto; justificar independencia y que es un plano: 1 punto.
    \item Obtener $m=4$: 0,5 puntos; expresar $w$ como combinación lineal: 0,5 puntos.
    \item Comprobar la pertenencia o no pertenencia del vector dado: 0,75 puntos; concluir: 0,25 puntos.
\end{itemize}
Como $u\times v=(7,-5,-3)\ne0$, los vectores son independientes y
\[
    \operatorname{span}\{u,v\}=\{a(1,2,-1)+b(2,1,3):a,b\in\R\}
    =\{(x,y,z):7x-5y-3z=0\}.
\]
Es un plano por el origen. Para $w$, la condición es
$7+5-3m=0$, así que $m=4$ y $w=v-u$.
Para $x=(-1,2,2)$, $7(-1)-5(2)-3(2)=-23\ne0$;
por tanto, $x$ no pertenece al espacio generado.

\section*{5. Subespacio de polinomios (4 puntos)}
\noindent\textbf{Rúbrica.}
\begin{itemize}[itemsep=2pt]
    \item Verificar la pertenencia del polinomio cero: 1 punto.
    \item Demostrar cerradura bajo suma usando ambas restricciones: 1,5 puntos.
    \item Demostrar cerradura bajo producto por escalar usando ambas restricciones: 1 punto.
    \item Concluir que es subespacio: 0,5 puntos.
\end{itemize}
El polinomio cero satisface las dos condiciones. Si $p,q\in W$ y
$\lambda\in\R$, entonces
\[
    (p+q)(2)=0,\quad (p+q)'(0)=p'(0)+q'(0)=0,
\]
\[
    (\lambda p)(2)=\lambda p(2)=0,\quad
    (\lambda p)'(0)=\lambda p'(0)=0.
\]
Como estas operaciones permanecen en $\mathcal P_3(\R)$, $W$ es subespacio.
Se acepta también identificarlo como intersección de núcleos de las
aplicaciones lineales $p\mapsto p(2)$ y $p\mapsto p'(0)$.
\end{document}
